\section{Reproducing the verification}\label{app:reproduce}

Everything needed to repeat \cref{sec:verification} is in the repository
\begin{center}
\url{https://github.com/vincitamore/riemann-lean-account}
\end{center}
together with the source of this paper, the reading notes, the delegated reports, the reviews
of the replay program and the audit. The archive itself is not redistributed there; it is
fetched from the author's site and checked against the hash of \cref{sec:verification}.

\subsection*{Requirements}
The first build ran on a rented Linux machine with \buildthreads{} threads whose memory was
limited to 128\,GB, and that is what gives the build time quoted here; the scripts' comments
estimate 150\,GB of free disk.
The scripts in \lean{verify/} assume \lean{bash} with \lean{git},
\lean{curl}, \lean{zstd}, \lean{tar} and \lean{sha256sum}, and \lean{python3} for the scan;
\lean{setup.sh} installs \lean{elan} if it is absent.

\subsection*{Steps}
\begin{enumerate}
\item \lean{setup.sh} downloads the fourteen parts, concatenates them, checks the SHA-256,
  unpacks, installs the pinned toolchain and compiles Mathlib from source at the pinned commit
  (or, on request, fetches its build cache).
\item \lean{scan.py} is the static inspection of \cref{sec:verification}: it reads the
  downloaded file without unpacking it and needs nothing else.
\item \lean{toolchain.sh} compares the installed toolchain with the release archive on GitHub.
\item \lean{build.sh} runs \lean{lake build Main} with every log line time-stamped and records
  the exit code.
\item \lean{janitor.sh} must run beside the build. Lake writes, for every module, a file listing
  the module's whole import closure%
\ifsecondrun
, \setupfilemb{}\,MB on average here; for this tree those
  files would total about \setupallgb{}\,GB. The script deletes each one once its module's
  compiled file exists, which Lake tolerates. With it the disk in use peaked at
  \diskpeak{}\,GB.
\else
, and by the script's estimate those files would for this tree exceed the disk by a wide
  margin. The script deletes each
  one once its module's compiled file exists, which Lake tolerates.
\fi
\item \lean{summary.sh} extracts from the logs and the built tree the facts of the build that
  this paper quotes.
\item \lean{check.sh} prints the statement of the main theorem, the module in which Mathlib's
  statement is declared, and the axioms of the five theorems.
\item \lean{cut.sh} builds and runs the program \lean{cut/CutMain.lean}, which computes the
  dependency cone of \cref{sec:cone} and writes the list of every declaration of the archive
  inside it and outside it.
\item \lean{replay.sh} builds and runs the program \lean{replay/ReplayMain.lean}, the kernel
  replay of \cref{sec:replay}. The number of threads is set with \lean{LEAN_NUM_THREADS}.
\item \lean{hostile.sh} builds the modules of \lean{hostile/} and runs Lean's checker and the
  replay program on each (\cref{tab:hostile}).
\item \lean{stock.sh} runs \lean{leanchecker --fresh} and the replay program on Mathlib, and
  \lean{sample.sh} runs \lean{leanchecker} on a sample of the archive's modules
  (\cref{sec:checker}).
\item \lean{independent.sh} fetches \lean{lean4lean} at a fixed commit, builds it with the
  archive's toolchain, and runs it on the modules of \lean{hostile/}, on the closure of the
  statement and on the same sample.
\end{enumerate}
The outputs of our runs are in \lean{verify/results/}, with an index that says what
wrote each file\ifsecondrun; those of the first run are in \lean{verify/results/run1/}\fi.
\ifsecondrun\else
In this version the files for the archive are those of the first build, and not all of them
were written by the scripts in their present form. The summary of the build was assembled by
hand from the build log, which was not kept, in a shorter form than \lean{summary.sh} now
writes. The cone was computed by the program as published with version~1.1, whose traversal
is the present one; the present program prints more. The replay was made by the first form
of the replay program, which is published beside its output. Steps of the list above that
concern a sample of the archive's modules belong to the second run and are not reported
here.
\fi
Every number in
this paper that was measured is read from those files by \lean{paper/tools/numbers.py}, which
also checks that they agree with one another, apart from three that no file holds and that
are typed with their provenance in \lean{version.tex}; and \lean{paper/tools/check_names.py}
checks every Lean name cited in the paper against the lists of declarations inside and
outside the cone.

\subsection*{Cost}
\ifsecondrun
The second run took \mathlibminutes{} minutes to compile Mathlib and \buildhours{} hours to
build the archive on \buildthreads{} threads; the replay then took \replayminutes{} minutes.
On rented hardware the two runs together cost about \totalcost{} US dollars.
\else
The build took \buildhours{} hours on \buildthreads{} threads and the replay \replayminutes{}
minutes. The provider billed \firstruncost{} US dollars for that machine, from its start to
its end.
\fi
